\chapter{2002年京蒙皖卷（春理）}

\section{单选题}
\begin{problem}
不等式组 \(\begin{cases}
x^{2}-1<0,\\
x^{2}-3x<0
\end{cases}\) 的解集为（\quad）
\choices
    {\(\{x\mid -1<x<1\}\)}
    {\(\{x\mid 0<x<3\}\)}
    {\(\{x\mid 0<x<1\}\)}
    {\(\{x\mid -1<x<3\}\)}
\end{problem}

\begin{answer}
C
\end{answer}

\begin{solution}
由 \(x^{2}-1<0\)，得 \((x-1)(x+1)<0\)，所以 \(-1<x<1\). 由 \(x^{2}-3x<0\)，得 \(x(x-3)<0\)，所以 \(0<x<3\). 两个不等式同时成立时，\(x\in(-1,1)\cap(0,3)=(0,1)\)，故选 C.
\end{solution}

\begin{problem}
已知三条直线 \(m\)，\(n\)，\(l\)，三个平面 \(\alpha\)，\(\beta\)，\(\gamma\)，下列四个命题中，正确的是（\quad）
\choices
    {\(\begin{cases}
\alpha\perp\gamma,\\
\beta\perp\gamma
\end{cases}\Rightarrow\alpha\parallel\beta\)}
    {\(\begin{cases}
m\parallel\beta,\\
l\perp m
\end{cases}\Rightarrow l\perp\beta\)}
    {\(\begin{cases}
m\parallel\gamma,\\
n\parallel\gamma
\end{cases}\Rightarrow m\parallel n\)}
    {\(\begin{cases}
m\perp\gamma,\\
n\perp\gamma
\end{cases}\Rightarrow m\parallel n\)}
\end{problem}

\begin{answer}
D
\end{answer}

\begin{solution}
两个平面都垂直于同一个平面时，它们可能相交，故命题 A 不一定成立. 直线 \(m\parallel\beta\) 且 \(l\perp m\) 时，\(l\) 不一定垂直于平面 \(\beta\)，故命题 B 不一定成立. 两条直线分别平行于同一个平面时，它们可能相交或异面，故命题 C 不一定成立. 垂直于同一个平面的两条直线都与该平面的法线方向平行，因此互相平行，命题 D 正确，故选 D.
\end{solution}

\begin{problem}
已知椭圆的焦点是 \(F_{1}\)，\(F_{2}\)，\(P\) 是椭圆上的一个动点. 如果延长 \(F_{1}P\) 到 \(Q\)，使得 \(\lvert PQ\rvert=\lvert PF_{2}\rvert\)，那么动点 \(Q\) 的轨迹是（\quad）
\choices
    {圆}
    {椭圆}
    {双曲线的一支}
    {抛物线}

\end{problem}

\begin{answer}
A
\end{answer}

\begin{solution}
设椭圆的长半轴为 \(a\). 对椭圆上任意点 \(P\)，由椭圆定义有 \(\lvert F_{1}P\rvert+\lvert PF_{2}\rvert=2a\). 由于 \(Q\) 在从 \(F_{1}\) 经过 \(P\) 的射线上，且 \(\lvert PQ\rvert=\lvert PF_{2}\rvert\)，所以
\(\lvert F_{1}Q\rvert=\lvert F_{1}P\rvert+\lvert PQ\rvert=\lvert F_{1}P\rvert+\lvert PF_{2}\rvert=2a\). 因而 \(Q\) 到定点 \(F_{1}\) 的距离恒为 \(2a\). 反过来，椭圆上各点从焦点 \(F_{1}\) 出发的射线覆盖所有方向，所以 \(Q\) 遍历以 \(F_{1}\) 为圆心、\(2a\) 为半径的圆，故选 A.
\end{solution}

\begin{problem}
如果 \(\theta\in\left(\dfrac{\pi}{2},\pi\right)\)，那么复数 \((1+\i)(\cos\theta+\i\sin\theta)\) 的辐角的主值是（\quad）
\choices
    {\(\theta+\dfrac{9\pi}{4}\)}
    {\(\theta+\dfrac{\pi}{4}\)}
    {\(\theta-\dfrac{\pi}{4}\)}
    {\(\theta+\dfrac{7\pi}{4}\)}
\end{problem}

\begin{answer}
B
\end{answer}

\begin{solution}
有 \(1+\i=\sqrt{2}\left(\cos\dfrac{\pi}{4}+\i\sin\dfrac{\pi}{4}\right)\)，故
\(\left(1+\i\right)\left(\cos\theta+\i\sin\theta\right)=\sqrt{2}\left(\cos\left(\theta+\dfrac{\pi}{4}\right)+\i\sin\left(\theta+\dfrac{\pi}{4}\right)\right)\). 由 \(\dfrac{\pi}{2}<\theta<\pi\)，得 \(\dfrac{3\pi}{4}<\theta+\dfrac{\pi}{4}<\dfrac{5\pi}{4}<2\pi\). 因此在辐角主值取 \([0,2\pi)\) 的约定下，主值就是 \(\theta+\dfrac{\pi}{4}\)，故选 B.
\end{solution}

\begin{problem}
若角 \(\alpha\) 满足条件 \(\sin2\alpha<0\)，\(\cos\alpha-\sin\alpha<0\)，则 \(\alpha\) 在（\quad）
\choices
    {第一象限}
    {第二象限}
    {第三象限}
    {第四象限}
\end{problem}

\begin{answer}
B
\end{answer}

\begin{solution}
在一个周期内，\(\sin2\alpha<0\) 时，\(\alpha\) 只能位于第二或第四象限. 在第二象限，\(\cos\alpha<0\)，\(\sin\alpha>0\)，所以 \(\cos\alpha-\sin\alpha<0\). 在第四象限，\(\cos\alpha>0\)，\(\sin\alpha<0\)，所以 \(\cos\alpha-\sin\alpha>0\)，不符合条件. 严格不等式又排除了坐标轴，故 \(\alpha\) 在第二象限，选 B.
\end{solution}

\begin{problem}
从 \(6\) 名志愿者中选出 \(4\) 人分别从事翻译、导游、导购、保洁 \(4\) 项不同工作. 若其中甲、乙两名支援者都不能从事翻译工作，则选派方案共有（\quad）
\choices
    {\(280\) 种}
    {\(240\) 种}
    {\(180\) 种}
    {\(96\) 种}

\end{problem}

\begin{answer}
B
\end{answer}

\begin{solution}
甲、乙不能从事翻译工作，所以翻译人员只能从其余 \(4\) 人中选出，方法数为 \(4\). 确定翻译人员后，从剩下的 \(5\) 人中选出 \(3\) 人并分别安排导游、导购、保洁，方法数为 \(\mathrm A_{5}^{3}=5\times4\times3=60\). 因此选派方案数为 \(4\times60=240\)，故选 B.
\end{solution}

\begin{problem}
在 \(\triangle ABC\) 中，\(AB=2\)，\(BC=1.5\)，\(\angle ABC=120^{\circ}\). 若使三角形绕直线 \(BC\) 旋转一周，则所形成的几何体的体积是（\quad）
\choices
    {\(\dfrac{3}{2}\pi\)}
    {\(\dfrac{5}{2}\pi\)}
    {\(\dfrac{7}{2}\pi\)}
    {\(\dfrac{9}{2}\pi\)}

\end{problem}

\begin{answer}
A
\end{answer}

\begin{solution}
从 \(A\) 向直线 \(BC\) 作垂线，垂足为 \(H\). 在直角三角形 \(ABH\) 中，
\(AH=AB\sin120^{\circ}=2\cdot\dfrac{\sqrt{3}}{2}=\sqrt{3}\)，且 \(BH=\lvert AB\cos120^{\circ}\rvert=1\). 因为 \(\angle ABC=120^{\circ}\)，垂足 \(H\) 在 \(B\) 的外侧，故 \(CH=BC+BH=\dfrac{3}{2}+1=\dfrac{5}{2}\).

旋转后，三角形所扫出的几何体等于以 \(C\) 为顶点、底面半径 \(AH\)、高 \(CH\) 的圆锥，减去以 \(B\) 为顶点、底面半径 \(AH\)、高 \(BH\) 的圆锥. 所以体积为
\(V=\dfrac{1}{3}\pi AH^{2}(CH-BH)=\dfrac{1}{3}\pi\cdot3\cdot\left(\dfrac{5}{2}-1\right)=\dfrac{3}{2}\pi\)，故选 A.
\end{solution}

\begin{problem}
圆 \(2x^{2}+2y^{2}=1\) 与直线 \(x\sin\theta+y-1=0\)（\(\theta\in\R\)，\(\theta\ne\dfrac{\pi}{2}+k\pi\)，\(k\in\Z\)）的位置关系是（\quad）
\choices
    {相交}
    {相切}
    {相离或相切}
    {不能确定}
\end{problem}

\begin{answer}
C
\end{answer}

\begin{solution}
圆 \(2x^{2}+2y^{2}=1\) 的圆心为原点，半径为 \(\dfrac{1}{\sqrt{2}}\). 直线 \(x\sin\theta+y-1=0\) 到原点的距离为
\(d=\dfrac{1}{\sqrt{1+\sin^{2}\theta}}\). 题设 \(\theta\ne\dfrac{\pi}{2}+k\pi\) 保证 \(\lvert\sin\theta\rvert<1\)，所以 \(d>\dfrac{1}{\sqrt{2}}\). 因而圆与直线相离. 选项中包含“相离”的只有“相离或相切”，故选 C.
\end{solution}

\begin{problem}
在极坐标系中，如果一个圆的方程 \(\rho=4\cos\theta+6\sin\theta\)，那么过圆心且与极轴平行的直线方程是（\quad）
\choices
    {\(\rho\sin\theta=3\)}
    {\(\rho\sin\theta=-3\)}
    {\(\rho\cos\theta=2\)}
    {\(\rho\cos\theta=-2\)}
\end{problem}

\begin{answer}
A
\end{answer}

\begin{solution}
由极坐标与直角坐标的关系 \(x=\rho\cos\theta\)，\(y=\rho\sin\theta\)，原方程化为 \(\rho^{2}=4\rho\cos\theta+6\rho\sin\theta\)，即 \(x^{2}+y^{2}=4x+6y\). 配方得 \((x-2)^{2}+(y-3)^{2}=13\)，所以圆心为 \((2,3)\). 过圆心且与极轴平行的直线是 \(y=3\)，其极坐标方程为 \(\rho\sin\theta=3\)，故选 A.
\end{solution}

\begin{problem}
对于二项式 \(\left(\dfrac1x+x^{3}\right)^{n}\) 的展开式（\(n\in\N^{*}\)），四位同学作出了四种判断：
\begin{circlelist}
\item 存在 \(n\in\N^{*}\)，展开式中有常数项；
\item 对任意 \(n\in\N^{*}\)，展开式中没有常数项；
\item 对任意 \(n\in\N^{*}\)，展开式中没有 \(x\) 的一次项；
\item 存在 \(n\in\N^{*}\)，展开式中有 \(x\) 的一次项.
\end{circlelist}
上述判断中正确的是（\quad）
\choices
    {\circled{1}与\circled{3}}
    {\circled{2}与\circled{3}}
    {\circled{1}与\circled{4}}
    {\circled{2}与\circled{4}}
\end{problem}

\begin{answer}
C
\end{answer}

\begin{solution}
展开式通项为 \(\mathrm C_{n}^{k}\left(\dfrac{1}{x}\right)^{n-k}\left(x^{3}\right)^{k}=\mathrm C_{n}^{k}x^{4k-n}\)，其中 \(k=0\)，\(1\)，\(\ldots\)，\(n\).

常数项要求 \(4k-n=0\)，取 \(n=4\)，\(k=1\) 即可满足，因此判断 \(\circled{1}\) 正确，判断 \(\circled{2}\) 错误. 一次项要求 \(4k-n=1\)，取 \(k=1\)，\(n=3\) 即可满足，因此判断 \(\circled{4}\) 正确，判断 \(\circled{3}\) 错误. 所以正确的判断是 \(\circled{1}\) 与 \(\circled{4}\)，故选 C.
\end{solution}

\begin{problem}
若一个等差数列前 3 项的和为 34，最后 3 项的和为 146，且所有项的和为 390，则这个数列有（\quad）
\choices
    {\(13\) 项}
    {\(12\) 项}
    {\(11\) 项}
    {\(10\) 项}

\end{problem}

\begin{answer}
A
\end{answer}

\begin{solution}
设等差数列为 \(a_{1}\)，\(a_{2}\)，\(\ldots\)，\(a_{n}\)，公差为 \(d\). 前三项的和为 \(3a_{2}=34\)，最后三项的和为 \(3a_{n-1}=146\)，所以 \(a_{2}=\dfrac{34}{3}\)，\(a_{n-1}=\dfrac{146}{3}\). 又 \(a_{1}=a_{2}-d\)，\(a_{n}=a_{n-1}+d\)，从而 \(a_{1}+a_{n}=a_{2}+a_{n-1}=60\). 因此总和为 \(390=\dfrac{n}{2}(a_{1}+a_{n})=30n\)，解得 \(n=13\)，故选 A.
\end{solution}

\begin{problem}
用一张钢板制作一个容积为 \(4\mathrm{~m}^{3}\) 的无盖长方体水箱. 可用的长方形钢板有四种不同的规格（长\(\times\)宽的尺寸如各选项所示，单位均为 \(\mathrm{m}\)），若既要够用，又要所剩最少，则应选钢板的规格是（\quad）
\choices
    {\(2\times5\)}
    {\(2\times5.5\)}
    {\(2\times6.1\)}
    {\(3\times5\)}
\end{problem}

\begin{answer}
C
\end{answer}

\begin{solution}
设水箱底面长、宽、高分别为 \(x\)，\(y\)，\(z\)，则 \(xyz=4\). 无盖水箱所需钢板面积为 \(S=xy+2xz+2yz\). 由基本不等式，
\(S\geqslant3\sqrt[3]{(xy)(2xz)(2yz)}=3\sqrt[3]{4x^{2}y^{2}z^{2}}=12\). 因而钢板面积至少为 \(12\mathrm{~m}^{2}\)，规格 \(2\times5\) 和 \(2\times5.5\) 的面积分别只有 \(10\mathrm{~m}^{2}\) 和 \(11\mathrm{~m}^{2}\)，不可能够用.

当 \(x=y=2\)，\(z=1\) 时，容积为 \(4\mathrm{~m}^{3}\)，所需五块钢板的尺寸为一块 \(2\times2\) 的底板和四块 \(2\times1\) 的侧板. 规格 \(2\times6.1\) 可沿长边排下这五块板，总长度为 \(2+4\times1=6\)，剩余面积为 \(2\times0.1=0.2\mathrm{~m}^{2}\). 规格 \(3\times5\) 也能排下这些板：先在一端放置 \(2\times2\) 的底板，余下的 \(2\times3\) 区域放置三块 \(2\times1\) 的侧板，再在剩下的 \(1\times2\) 区域放置一块旋转 \(90^\circ\) 后的侧板. 但其面积为 \(15\mathrm{~m}^{2}\)，剩余面积为 \(3\mathrm{~m}^{2}\). 因此在够用的规格中，\(2\times6.1\) 所剩最少，故选 C.
\end{solution}

\section{填空题}

\begin{problem}
若双曲线 \(\dfrac{x^{2}}{4}-\dfrac{y^{2}}{m}=1\) 的渐近线方程为 \(y=\pm\dfrac{\sqrt{3}}{2}x\)，则双曲线的焦点坐标是\fillinblank{}.
\end{problem}

\begin{answer}
\(\left(-\sqrt{7},0\right)\)，\(\left(\sqrt{7},0\right)\)
\end{answer}

\begin{solution}
双曲线的标准形式为 \(\dfrac{x^{2}}{a^{2}}-\dfrac{y^{2}}{b^{2}}=1\)，所以 \(a^{2}=4\)，\(a=2\)，且 \(b^{2}=m\). 其渐近线为 \(y=\pm\dfrac{b}{a}x\). 由题设 \(\dfrac{b}{2}=\dfrac{\sqrt{3}}{2}\)，得 \(b=\sqrt{3}\)，从而 \(m=b^{2}=3\). 双曲线的半焦距满足 \(c^{2}=a^{2}+b^{2}=4+3=7\)，故两个焦点坐标为 \(\left(-\sqrt{7},0\right)\)，\(\left(\sqrt{7},0\right)\).
\end{solution}

\begin{problem}
如果 \(\cos\theta=-\dfrac{12}{13}\)，\(\theta\in\left(\pi,\dfrac{3\pi}{2}\right)\)，那么 \(\cos\left(\theta+\dfrac{\pi}{4}\right)\) 的值等于\fillinblank{}.
\end{problem}

\begin{answer}
\(-\dfrac{7\sqrt{2}}{26}\)
\end{answer}

\begin{solution}
因为 \(\theta\in\left(\pi,\dfrac{3\pi}{2}\right)\)，所以 \(\sin\theta<0\). 由 \(\cos\theta=-\dfrac{12}{13}\)，得
\(\sin\theta=-\sqrt{1-\cos^{2}\theta}=-\sqrt{1-\dfrac{144}{169}}=-\dfrac{5}{13}\). 于是由余弦和角公式，
\(\cos\left(\theta+\dfrac{\pi}{4}\right)=\cos\theta\cos\dfrac{\pi}{4}-\sin\theta\sin\dfrac{\pi}{4}=\left(-\dfrac{12}{13}+\dfrac{5}{13}\right)\dfrac{\sqrt{2}}{2}=-\dfrac{7\sqrt{2}}{26}\).
\end{solution}

\begin{problem}
正方形 \(ABCD\) 的边长是 \(2\)，\(E\)，\(F\) 分别是 \(AB\) 和 \(CD\) 的中点，将正方形沿 \(EF\) 折成直二面角（如图所示）. \(M\) 为矩形 \(AEFD\) 内一点，如果 \(\angle MBE=\angle MBC\)，\(MB\) 和平面 \(BCF\) 所成角的正切值为 \(\dfrac{1}{2}\)，那么点 \(M\) 到直线 \(EF\) 的距离为\fillinblank{}.
\bitmapfigure[max width=0.52\linewidth]{img/2002/jing_meng_wan_science_spring/q15_fig1.png}
\end{problem}

\begin{answer}
\(\dfrac{\sqrt{2}}{2}\)
\end{answer}

\begin{solution}
折叠后建立空间直角坐标系，令折痕 \(EF\) 为 \(x\) 轴，取 \(E=(0,0,0)\)，\(F=(2,0,0)\)，使平面 \(BCF\) 为 \(z=0\)，并令 \(B=(0,1,0)\)，\(C=(2,1,0)\). 由于 \(M\) 在矩形 \(AEFD\) 内，可设 \(M=(t,0,h)\)，其中 \(0<t<2\)，\(0<h<1\).

于是 \(\overrightarrow{BM}=(t,-1,h)\)，\(\overrightarrow{BE}=(0,-1,0)\)，\(\overrightarrow{BC}=(2,0,0)\). 由题设两角相等，且两角均为锐角，比较余弦得
\(\dfrac{\overrightarrow{BM}\cdot\overrightarrow{BE}}{\lvert BM\rvert\lvert BE\rvert}=\dfrac{1}{\sqrt{t^{2}+1+h^{2}}}\)，\(\dfrac{\overrightarrow{BM}\cdot\overrightarrow{BC}}{\lvert BM\rvert\lvert BC\rvert}=\dfrac{t}{\sqrt{t^{2}+1+h^{2}}}\). 因而 \(t=1\).

直线 \(MB\) 在平面 \(BCF\) 上的射影长度为 \(\sqrt{t^{2}+1}=\sqrt{2}\). 设 \(MB\) 与平面 \(BCF\) 所成的角为 \(\varphi\)，则 \(\tan\varphi=\dfrac{h}{\sqrt{2}}=\dfrac{1}{2}\)，所以 \(h=\dfrac{\sqrt{2}}{2}\). 此时 \(M=(1,0,h)\)，而 \(EF\) 是 \(x\) 轴，故点 \(M\) 到直线 \(EF\) 的距离为 \(\sqrt{0^{2}+h^{2}}=h=\dfrac{\sqrt{2}}{2}\).
\end{solution}

\begin{problem}
对于任意两个复数 \(z_{1}=x_{1}+y_{1}\i\)，\(z_{2}=x_{2}+y_{2}\i\)（\(x_{1}\in\R\)，\(y_{1}\in\R\)，\(x_{2}\in\R\)，\(y_{2}\in\R\)），定义运算 \(\odot\) 为：\(z_{1}\odot z_{2}=x_{1}x_{2}+y_{1}y_{2}\). 设非零复数 \(w_{1}\)，\(w_{2}\) 在复平面内对应的点分别为 \(P_{1}\)，\(P_{2}\)，点 \(O\) 为坐标原点. 如果 \(w_{1}\odot w_{2}=0\)，那么在 \(\triangle P_{1}OP_{2}\) 中，\(\angle P_{1}OP_{2}\) 的大小为\fillinblank{}.
   \end{problem}

\begin{answer}
\(\dfrac{\pi}{2}\)
\end{answer}

\begin{solution}
复数 \(w_{1}\)，\(w_{2}\) 对应的向量分别为 \(\overrightarrow{OP_{1}}=(x_{1},y_{1})\)，\(\overrightarrow{OP_{2}}=(x_{2},y_{2})\). 由新定义，
\(\overrightarrow{OP_{1}}\cdot\overrightarrow{OP_{2}}=x_{1}x_{2}+y_{1}y_{2}=w_{1}\odot w_{2}=0\). 由于 \(w_{1}\)，\(w_{2}\) 均非零，两个向量均非零；非零向量内积为零等价于两向量垂直. 因此 \(\angle P_{1}OP_{2}=\dfrac{\pi}{2}\).
\end{solution}

\section{解答题}

\begin{problem}
在 \(\triangle ABC\) 中，已知 \(A\)，\(B\)，\(C\) 成等差数列，求 \(\tan\dfrac{A}{2}+\tan\dfrac{C}{2}+\sqrt{3}\tan\dfrac{A}{2}\tan\dfrac{C}{2}\) 的值.

\begin{solution}
因为 \(A\)，\(B\)，\(C\) 成等差数列，所以 \(A+C=2B\). 又因为 \(A+B+C=\pi\)，所以 \(3B=\pi\)，从而 \(B=\dfrac{\pi}{3}\)，且 \(\dfrac{A+C}{2}=\dfrac{\pi}{3}\). 设 \(x=\tan\dfrac{A}{2}\)，\(y=\tan\dfrac{C}{2}\). 由于 \(A>0\)，\(C>0\)，且 \(\dfrac{A}{2}+\dfrac{C}{2}=\dfrac{\pi}{3}<\dfrac{\pi}{2}\)，有 \(\tan\dfrac{A+C}{2}=\dfrac{x+y}{1-xy}=\sqrt{3}\). 因此 \(x+y=\sqrt{3}(1-xy)\)，即 \(x+y+\sqrt{3}xy=\sqrt{3}\). 所以所求值为 \(\sqrt{3}\).
\end{solution}
\end{problem}

\begin{problem}
已知函数 \(f(x)\) 是偶函数，而且在 \((0,+\infty)\) 上是增函数，判断 \(f(x)\) 在 \(\left(-\infty,0\right)\) 上是增函数还是减函数，并证明你的判断.

\begin{solution}
在 \(\left(-\infty,0\right)\) 上是减函数. 任取 \(x_{1}\)，\(x_{2}\in\left(-\infty,0\right)\)，且 \(x_{1}<x_{2}<0\)，则 \(-x_{1}>-x_{2}>0\). 由 \(f(x)\) 在 \((0,+\infty)\) 上是增函数，得 \(f(-x_{1})>f(-x_{2})\). 又因为 \(f(x)\) 是偶函数，所以 \(f(x_{1})=f(-x_{1})\)，\(f(x_{2})=f(-x_{2})\)，从而 \(f(x_{1})>f(x_{2})\). 因此，当 \(x_{1}<x_{2}<0\) 时总有 \(f(x_{1})>f(x_{2})\)，故 \(f(x)\) 在 \(\left(-\infty,0\right)\) 上是减函数.
\end{solution}
\end{problem}

\begin{problem}
在三棱锥 \(S-ABC\) 中，如图，\(\angle SAB=\angle SAC=\angle ACB=90^{\circ}\)，\(AC=2\)，\(BC=\sqrt{13}\)，\(SB=\sqrt{29}\).
\begin{enumerate}
\item 证明：\(SC\perp BC\)；
\item 求侧面 \(SBC\) 与底面 \(ABC\) 所成的二面角大小；
\item 求异面直线 \(SC\) 与 \(AB\) 所成的角的大小（用反三角函数表示）.
\end{enumerate}
\bitmapfigure[max width=0.42\linewidth]{img/2002/jing_meng_wan_science_spring/q19_fig1.png}

\begin{solution}
\begin{enumerate}
\item 由 \(\angle SAB=\angle SAC=90^{\circ}\)，且直线 \(AB\)，\(AC\) 在平面 \(ABC\) 内相交于 \(A\)，可知 \(SA\perp\text{平面 }ABC\). 取直角坐标系，使 \(A=(0,0,0)\)，\(C=(2,0,0)\)，\(B=(2,\sqrt{13},0)\)，则由 \(SB=\sqrt{29}\) 得 \(S=(0,0,h)\)，且 \(4+13+h^{2}=29\)，所以 \(h=2\sqrt{3}\). 因此 \(\overrightarrow{SC}=(2,0,-2\sqrt{3})\)，\(\overrightarrow{BC}=(0,-\sqrt{13},0)\)，从而 \(\overrightarrow{SC}\cdot\overrightarrow{BC}=0\)，故 \(SC\perp BC\).
\item 由第（1）问，直线 \(SC\) 垂直于棱 \(BC\)，而直线 \(AC\) 也垂直于 \(BC\)，所以侧面 \(SBC\) 与底面 \(ABC\) 所成的二面角等于 \(\angle ACS\). 又 \(\overrightarrow{CA}=(-2,0,0)\)，\(\overrightarrow{CS}=(-2,0,2\sqrt{3})\)，故 \(\lvert CA\rvert=2\)，\(\lvert CS\rvert=4\)，且 \(\overrightarrow{CA}\cdot\overrightarrow{CS}=4\). 因而 \(\cos\angle ACS=\dfrac{4}{2\times4}=\dfrac12\)，所以二面角大小为 \(\dfrac{\pi}{3}\).
\item 直线 \(SC\) 与 \(AB\) 的方向向量可分别取 \(\overrightarrow{CS}=(-2,0,2\sqrt{3})\) 和 \(\overrightarrow{AB}=(2,\sqrt{13},0)\). 两向量的数量积为 \(-4\)，长度分别为 \(4\) 和 \(\sqrt{17}\). 设异面直线所成的锐角为 \(\varphi\)，则 \(\cos\varphi=\dfrac{\lvert-4\rvert}{4\sqrt{17}}=\dfrac{1}{\sqrt{17}}\)，所以 \(\varphi=\arccos\dfrac{1}{\sqrt{17}}=\arctan4\).
\end{enumerate}
\end{solution}
\end{problem}

\begin{problem}
假设 \(A\) 型进口车关税税率在 \(2001\text{ 年}\)是 \(100\%\)，在 \(2006\text{ 年}\)是 \(25\%\)，\(2001\text{ 年}\) \(A\) 型进口车每辆价格为 \(64\text{ 万元}\)（其中含 \(32\text{ 万元}\)关税税款）.
\begin{enumerate}
\item 已知与 \(A\) 型车性能相近的 \(B\) 型国产车，\(2001\text{ 年}\)每辆价格为 \(46\text{ 万元}\)，若 \(A\) 型车的价格只受关税降低的影响，为了保证 \(2006\text{ 年}\) \(B\) 型车的价格不高于 \(A\) 型车价格的 \(90\%\)，\(B\) 型车价格要逐年降低，问平均每年至少下降多少万元？
\item 某人在 \(2001\text{ 年}\)将 \(33\text{ 万元}\)存入银行，假设银行扣利息税后的年利率为 \(1.8\%\)（\(5\text{ 年}\)内不变），且每年按复利计算（上一年的利息计入第二年的本金），那么五年到期时这笔钱连本带息是否一定够买按（1）中所述降价后的 \(B\) 型车一辆？
\end{enumerate}

\begin{solution}
\begin{enumerate}
\item 设 \(A\) 型车的完税前价格为 \(p\) 万元. 由 \(2001\text{ 年}\)的关税税率为 \(100\%\)，且关税为 \(32\text{ 万元}\)，得 \(p=32\). 因而 \(2006\text{ 年}\)的 \(A\) 型车价格为 \(32(1+25\%)=40\) 万元. 为使 \(B\) 型车价格不高于该价格的 \(90\%\)，其 \(2006\text{ 年}\)价格至多为 \(40\times90\%=36\) 万元. 由 \(2001\text{ 年}\)的 \(46\) 万元降至至多 \(36\) 万元，五年间至少下降 \(46-36=10\) 万元. 因此平均每年至少下降 \(\dfrac{10}{2006-2001}=2\) 万元.
\item 按每年至少下降 \(2\) 万元计算，\(2006\text{ 年}\)的 \(B\) 型车价格至多为 \(36\) 万元. 五年后存款本息为 \(33(1+1.8\%)^{5}=33(1.018)^{5}\approx36.08\) 万元，而 \(36.08>36\). 所以五年到期时这笔钱一定够买按上述降价后的 \(B\) 型车一辆.
\end{enumerate}
\end{solution}
\end{problem}

\begin{problem}
已知点的序列 \(A_n(x_n,0)\)，\(n\in\N\)，其中 \(x_1=0\)，\(x_2=a\)（\(a<0\)），\(A_3\) 是线段 \(A_1A_2\) 的中点，\(A_4\) 是线段 \(A_2A_3\) 的中点，\(\cdots\)，\(A_n\) 是线段 \(A_{n-2}A_{n-1}\) 的中点，\(\cdots\).
\begin{enumerate}
\item 写出 \(x_n\) 与 \(x_{n-1}\)，\(x_{n-2}\) 之间的关系式（\(n\geqslant3\)）；
\item 设 \(a_n=x_{n+1}-x_n\)，计算 \(a_1\)，\(a_2\)，\(a_3\)，由此推测数列 \(\{a_n\}\) 的通项公式，并加以证明；
\item 求 \(\displaystyle\lim_{n\to\infty}x_n\).
\end{enumerate}

\begin{solution}
\begin{enumerate}
\item 因为 \(A_n\) 是线段 \(A_{n-2}A_{n-1}\) 的中点，所以 \(x_n=\dfrac{x_{n-2}+x_{n-1}}{2}\)，即 \(2x_n=x_{n-1}+x_{n-2}\)（\(n\geqslant3\)）.
\item 由 \(x_1=0\)，\(x_2=a\)，得 \(x_3=\dfrac{x_1+x_2}{2}=\dfrac{a}{2}\)，\(x_4=\dfrac{x_2+x_3}{2}=\dfrac{3a}{4}\). 因而 \(a_1=x_2-x_1=a\)，\(a_2=x_3-x_2=-\dfrac{a}{2}\)，\(a_3=x_4-x_3=\dfrac{a}{4}\). 对任意 \(n\geqslant1\)，由递推关系有 \(a_{n+1}=x_{n+2}-x_{n+1}=\dfrac{x_n+x_{n+1}}{2}-x_{n+1}=-\dfrac12a_n\). 所以 \(\{a_n\}\) 是首项为 \(a\)、公比为 \(-\dfrac12\) 的等比数列，从而 \(a_n=a\left(-\dfrac12\right)^{n-1}\).
\item 由 \(x_n=x_1+\displaystyle\sum_{k=1}^{n-1}a_k\)，得 \(x_n=a\displaystyle\sum_{k=0}^{n-2}\left(-\dfrac12\right)^k=\dfrac{2a}{3}\left(1-\left(-\dfrac12\right)^{n-1}\right)\). 因为 \(\left(-\dfrac12\right)^{n-1}\to0\)，所以 \(\displaystyle\lim_{n\to\infty}x_n=\dfrac{2a}{3}\).
\end{enumerate}
\end{solution}
\end{problem}

\begin{problem}
已知某椭圆的焦点是 \(F_1(-4,0)\)，\(F_2(4,0)\)，过点 \(F_2\)，并垂直于 \(x\) 轴的直线与椭圆的一个交点为 \(B\)，且 \(\lvert F_1B\rvert+\lvert F_2B\rvert=10\). 椭圆上不同的两点 \(A(x_1,y_1)\)，\(C(x_2,y_2)\) 满足条件：\(\lvert F_2A\rvert\)，\(\lvert F_2B\rvert\)，\(\lvert F_2C\rvert\) 成等差数列.
\begin{enumerate}
\item 求该椭圆的方程；
\item 求弦 \(AC\) 中点的横坐标；
\item 设弦 \(AC\) 的垂直平分线的方程为 \(y=kx+m\)，求 \(m\) 的取值范围.
\end{enumerate}

\begin{solution}
\begin{enumerate}
\item 设椭圆的长半轴为 \(a_{0}\)，短半轴为 \(b_{0}\)，半焦距为 \(c_{0}\). 由焦点坐标得 \(c_{0}=4\). 又因为 \(B\) 在椭圆上，故 \(\lvert F_{1}B\rvert+\lvert F_{2}B\rvert=2a_{0}=10\)，从而 \(a_{0}=5\)，且 \(b_{0}^{2}=a_{0}^{2}-c_{0}^{2}=25-16=9\). 所以椭圆方程为 \(\dfrac{x^{2}}{25}+\dfrac{y^{2}}{9}=1\).
\item 对椭圆上任意点 \(P(x,y)\)，记 \(r_{1}=\lvert F_{1}P\rvert\)，\(r_{2}=\lvert F_{2}P\rvert\). 由椭圆定义，\(r_{1}+r_{2}=10\)，又 \(r_{1}^{2}-r_{2}^{2}=(x+4)^{2}+y^{2}-\bigl((x-4)^{2}+y^{2}\bigr)=16x\)，所以 \(r_{1}-r_{2}=\dfrac{8}{5}x\)，从而 \(\lvert F_{2}P\rvert=r_{2}=5-\dfrac45x\). 由第（1）问，点 \(B\) 的横坐标为 \(4\)，故 \(\lvert F_{2}B\rvert=5-\dfrac45\times4=\dfrac95\). 设 \(M\) 为弦 \(AC\) 的中点. 由三项成等差数列，\(\lvert F_{2}A\rvert+\lvert F_{2}C\rvert=2\lvert F_{2}B\rvert=\dfrac{18}{5}\). 若 \(M=(x_{M},y_{M})\)，则 \(x_{M}=\dfrac{x_{A}+x_{C}}{2}=\dfrac{1}{2}\left[\dfrac54\left(5-\lvert F_{2}A\rvert\right)+\dfrac54\left(5-\lvert F_{2}C\rvert\right)\right]=4\). 所以弦 \(AC\) 中点的横坐标为 \(4\).
\item 记弦 \(AC\) 的中点为 \(M=(4,y_{M})\)，并设 \(A=(x_{A},y_{A})\)，\(C=(x_{C},y_{C})\). 由两点在椭圆上，作差得 \(\dfrac{x_{A}^{2}-x_{C}^{2}}{25}+\dfrac{y_{A}^{2}-y_{C}^{2}}{9}=0\). 又 \(x_{A}+x_{C}=8\)，\(y_{A}+y_{C}=2y_{M}\)，所以弦的方向向量 \((x_{A}-x_{C},y_{A}-y_{C})\) 满足 \(\dfrac{8}{25}(x_{A}-x_{C})+\dfrac{2y_{M}}{9}(y_{A}-y_{C})=0\). 因而向量 \(\left(\dfrac{8}{25},\dfrac{2y_{M}}{9}\right)\) 是弦的一个法向量，弦的垂直平分线斜率为 \(k=\dfrac{(2y_{M}/9)}{(8/25)}=\dfrac{25y_{M}}{36}\)，这个推导也包含 \(y_{M}=0\) 的情形. 该直线经过 \(M\)，故 \(m=y_{M}-4k=-\dfrac{16}{9}y_{M}\).

又因为 \(M\) 是椭圆内一个非边界点，且 \(x_{M}=4\)，所以 \(\dfrac{16}{25}+\dfrac{y_{M}^{2}}{9}<1\)，即 \(-\dfrac95<y_{M}<\dfrac95\). 反之，直线 \(x=4\) 上任意满足此不等式的点均在椭圆内部，均可作为某条非退化弦的中点，故取值范围没有遗漏. 由 \(m=-\dfrac{16}{9}y_{M}\)，得 \(-\dfrac{16}{5}<m<\dfrac{16}{5}\).
\end{enumerate}
\end{solution}
\end{problem}
